% ===== PRÉAMBULE CANONIQUE — à coller VERBATIM en tête de CHAQUE .tex =====
\documentclass[11pt,a4paper]{article}
\usepackage[utf8]{inputenc}\usepackage[T1]{fontenc}\usepackage{lmodern}
\usepackage[french]{babel}
\usepackage{amsmath,amssymb,mathtools}\usepackage{icomma}
\usepackage{siunitx}\sisetup{locale=FR}  % \qty{}{}, \si{}, \ang{}
\usepackage{xcolor}\usepackage{colortbl}
\usepackage{tikz}\usepackage{pgfplots}\pgfplotsset{compat=1.18}
\usepackage[siunitx,europeanresistors]{circuitikz} % circuits (élec.)
\usepackage[most]{tcolorbox}\usepackage{enumitem}\usepackage{array,booktabs}
\usepackage[a4paper,margin=2.2cm]{geometry}
\usetikzlibrary{arrows.meta,calc,angles,quotes,positioning,patterns,%
  backgrounds,shapes.geometric,decorations.pathmorphing,decorations.markings,intersections}
\definecolor{primary}{RGB}{222,49,99}\definecolor{primarydark}{RGB}{150,22,55}
\definecolor{defcol}{RGB}{0,114,178}\definecolor{methcol}{RGB}{52,92,148}
\definecolor{excol}{RGB}{0,135,90}\definecolor{attncol}{RGB}{200,40,55}
\tcbset{boxbase/.style={enhanced,breakable,boxrule=0.9pt,arc=2.5pt,
  left=3mm,right=3mm,top=2.3mm,bottom=2.3mm,fonttitle=\bfseries,coltitle=white}}
% --- boîtes TUTORIEL (noms EXACTS, ne pas renommer) ---
\newtcolorbox{cle}[1][]{boxbase,colback=defcol!6,colframe=defcol,
  title={\textsc{À retenir}\ifx\relax#1\relax\relax\else\textmd{ : \itshape #1}\fi}}
\newtcolorbox{methode}[1][]{boxbase,colback=methcol!7,colframe=methcol,sharp corners,
  title={\textsc{Méthode}\ifx\relax#1\relax\relax\else\textmd{ --- \itshape #1}\fi}}
\newtcolorbox{exemplebox}[1][]{boxbase,colback=excol!5,colframe=excol,
  title={\textsc{Exemple}\ifx\relax#1\relax\relax\else\textmd{ : \itshape #1}\fi}}
\newtcolorbox{piege}[1][]{boxbase,colback=attncol!6,colframe=attncol,
  title={\textsc{Piège}\ifx\relax#1\relax\relax\else\textmd{ : \itshape #1}\fi}}
\newtcolorbox{remarque}[1][]{boxbase,colback=black!4,colframe=black!45,
  title={\textsc{Remarque}\ifx\relax#1\relax\relax\else\textmd{ : \itshape #1}\fi}}
% --- boîtes EXERCICE / CORRIGÉ (noms EXACTS, auto-comptées) ---
\newtcolorbox[auto counter]{exo}[1][]{boxbase,colback=primary!4,colframe=primary,
  title={\textsc{Exercice}~\thetcbcounter\ifx\relax#1\relax\relax\else\textmd{ --- \itshape #1}\fi}}
\newtcolorbox[auto counter]{corrige}[1][]{boxbase,colback=excol!4,colframe=excol,
  title={\textsc{Corrigé}~\thetcbcounter\ifx\relax#1\relax\relax\else\textmd{ --- \itshape #1}\fi}}
\newcommand{\vv}[1]{\vec{#1}}\newcommand{\ds}{\displaystyle}
\newcommand{\R}{\mathbb{R}}\newcommand{\N}{\mathbb{N}}\newcommand{\Z}{\mathbb{Z}}
% (un \usepackage{fancyhdr}+en-tête est possible ; il est ignoré côté web)
% =========================================================================
\begin{document}
\begin{exo}[poids : application directe]
Un corps a une masse $m=\qty{3}{kg}$.
\begin{enumerate}[label=\alph*)]
  \item Calcule son poids avec $g=\qty{10}{N/kg}$.
  \item Recalcule-le avec la valeur plus précise $g=\qty{9.81}{N/kg}$.
\end{enumerate}
\end{exo}

\begin{exo}[convertir avant de calculer]
Une bille a une masse $m=\qty{250}{g}$. On prend $g=\qty{10}{N/kg}$.
\begin{enumerate}[label=\alph*)]
  \item Convertis la masse en kilogrammes.
  \item Calcule le poids $G$ de la bille.
\end{enumerate}
\end{exo}

\begin{exo}[réaction d'une table]
Un bloc immobile est posé sur une table horizontale. Son poids vaut $G=\qty{15}{N}$.
\begin{center}
\begin{tikzpicture}[>=Stealth,scale=1]
  \draw[black!55,line width=1pt] (-2,0)--(2,0);
  \foreach \x in {-1.7,-1.3,...,1.8}\draw[black!45](\x,0)--(\x-0.22,-0.28);
  \draw[fill=primary!18,draw=primary!70] (-0.6,0) rectangle (0.6,0.55);
  \fill (0,0.275) circle (1.2pt);
  \draw[->,line width=1.2pt,orange] (0,0.275)--(0,-1.0) node[below]{$\vv G$};
  \draw[->,line width=1.2pt,defcol,densely dashed] (0,0.9)--(0,1.7) node[above]{$\vv R\,?$};
\end{tikzpicture}
\end{center}
\begin{enumerate}[label=\alph*)]
  \item Quelles sont les deux forces qui s'exercent sur le bloc ?
  \item En écrivant l'équilibre de translation, donne l'intensité $R$ de la réaction de la table.
\end{enumerate}
\end{exo}

\begin{exo}[objet suspendu à un fil]
Un objet de masse $m=\qty{2}{kg}$ pend, immobile, au bout d'un fil vertical accroché au plafond
($g=\qty{10}{N/kg}$).
\begin{center}
\begin{tikzpicture}[>=Stealth,scale=1]
  \draw[black!55,line width=1pt] (-1.4,2)--(1.4,2);
  \foreach \x in {-1.2,-0.8,...,1.3}\draw[black!45](\x,2)--(\x-0.2,2.22);
  \draw[defcol,line width=1.1pt] (0,2)--(0,0.4);
  \draw[fill=primary!25,draw=primary!70] (-0.35,-0.2) rectangle (0.35,0.4);
  \node at (0,0.1){\scriptsize $m$};
  \draw[->,line width=1.2pt,orange] (0,-0.2)--(0,-1.2) node[below]{$\vv G$};
  \node[defcol] at (0.45,1.3){$\vv T$};
\end{tikzpicture}
\end{center}
\begin{enumerate}[label=\alph*)]
  \item Calcule le poids $G$ de l'objet.
  \item Déduis-en la tension $T$ du fil.
\end{enumerate}
\end{exo}

\begin{exo}[le même corps, deux astres]
Un astronaute emporte un objet de masse $m=\qty{5}{kg}$. On donne $g_{\text{Terre}}=\qty{10}{N/kg}$
et $g_{\text{Lune}}=\qty{1.6}{N/kg}$.
\begin{enumerate}[label=\alph*)]
  \item Calcule le poids de l'objet sur la Terre, puis sur la Lune.
  \item La \emph{masse} de l'objet change-t-elle d'un astre à l'autre ? Où le poids est-il le plus grand ?
\end{enumerate}
\end{exo}

\end{document}
